\input{../tex-inputs/latex-leseansicht-vorspann}

\section[Gustav Schwarzkopf an Arthur Schnitzler, 17. 7. 1907]{L04271 Gustav Schwarzkopf an Arthur Schnitzler, 17. 7. 1907}
\nopagebreak\mylabel{L04271v}
\rehead{ }\normalsize\beginnumbering\briefempfaengerindex{Schnitzler, Arthur@\textsc{Schnitzler, Arthur}!zzzSchwarzkopf, Gustav@\emph{von Gustav Schwarzkopf}!1907-07-171@{17. 7. 1907}|(be}
\toendnotes[C]{\smallbreak\pagebreak[2]}
\correspDesc{Versand  durch Gustav Schwarzkopf am 17. 7. 1907 in Wien
\newline{}Übermittlung  am 18. 7. 1907 in Wien
\newline{}Erhalt  durch Arthur Schnitzler im Zeitraum [18. 7. 1907 – 22. 7. 1907?] in Welsberg-Taisten}\toendnotes[C]{\smallbreak}
\Standort{CUL, Schnitzler, B 96a.}
\physDesc{Postkarte, 754 Zeichen
\newline{}Handschrift: schwarze Tinte, deutsche Kurrent
\newline{}Versand: Stempel: »\nobreak{}\oindex{I., Innere Stadt@\textbf{I., Innere Stadt}, \emph{Verwaltungsgebiet}|pwk}1/1 Wien 8, 18. VII. 07, 1\nobreak{}«.  }\toendnotes[C]{\smallbreak}\pstart{}{\pb}I. Tiefer Graben 23\oindex{Wien@\textbf{Wien}!I., Innere Stadt@\textbf{I., Innere Stadt}!Tiefer Graben 23@\textbf{Tiefer Graben 23}, \emph{Wohngebäude}|pw}.\pend{}\pstart{}II. St. Th. 6.\pend{}{\bigskip}\pstart{}Herrn \textsc{D\textsuperscript{r} Arthur
                     Schnitzler}\pend{}\pstart{}Bad \textsc{Waldbrunn}\oindex{Wildbad Waldbrunn@\textbf{Wildbad Waldbrunn}, \emph{Spa}|pw}\pend{}\pstart{}Station \textsc{Welsberg
                        (Tirol)}\oindex{Bahnhof Welsberg-Gsies@\textbf{Bahnhof Welsberg-Gsies}, \emph{Bahnhofsgebäude}|pw}\pend{}{\bigskip}\vspace{1em}
\pstart
           \raggedleft{}{\pb}17. VII 07\pend
           
\pstart{}Lieber Arthur!\pend\vspace{0.5em}
\pstart
           Ich war eine Woche und Graz\oindex{Graz@\textbf{Graz}, \emph{Verwaltungsgebiet}|pw}, meine Anweſenheit
               war doch notwendig. Emil\pwindex{Schwarzkopf, Emil 17.\,9.\,1851 Wien – 28.\,1.\,1928 ebd.@\textsc{Schwarzkopf, Emil} (17.\,9.\,1851 Wien – 28.\,1.\,1928 ebd.), \emph{Übersetzer, Komponist}|pw} iſt jetzt im Rekonvalescentenhaus\oindex{Rekonvaleszentenheim der barmherzigen Brüder@\textbf{Rekonvaleszentenheim der barmherzigen Brüder}, \emph{Wohngebäude}|pw} in Eggenberg\oindex{Eggenberg [Graz]@\textbf{Eggenberg [Graz]}, \emph{Ehemaliger Ort}|pw}. Mit diesem Grazer\oindex{Graz@\textbf{Graz}, \emph{Verwaltungsgebiet}|pw}-Ausflug dürfte meine So{\geminationm}erplänen
               wahrſcheinlich ein Ziel geſetzt{ }ſein – indeſſen, man{ }ſoll nichts
               verſchwören. –\pend
           
\pstart
           Paul M.\pwindex{Marx, Paul 21.\,7.\,1879 Wien – 30.\,10.\,1956 ebd.@\textsc{Marx, Paul} (21.\,7.\,1879 Wien – 30.\,10.\,1956 ebd.), \emph{Regisseur, Schauspieler}|pw} hat mir einen Abſchiedsbeſuch gemacht,
               hat er Sie aufgeſucht? Ich habe ihm noch zugeredet, in Waldbru{\geminationn}\oindex{Wildbad Waldbrunn@\textbf{Wildbad Waldbrunn}, \emph{Spa}|pw} Station zu machen. – Hoffentlich
               war das Wetter bei Ihnen nichts ganz{ }ſo erbärmlich wie hier. Max\pwindex{Schwarzkopf, Max 12.\,6.\,1857 Wien – 14.\,4.\,1928 ebd.@\textsc{Schwarzkopf, Max} (12.\,6.\,1857 Wien – 14.\,4.\,1928 ebd.), \emph{Rechtsanwalt}|pw} wartet{ }ſeit 5 Tagen mit gepacktem Ruckſack auf günſtiges
               Wetter zum Abreiſen; er will auf den Hochſchwab\oindex{Hochschwab@\textbf{Hochschwab}, \emph{Gebirge}|pw}. – Grüßen Sie Ihre Frau\pwindex{Schnitzler, Olga 17.\,1.\,1882 Wien – 13.\,1.\,1970 Lugano@\textsc{Schnitzler, Olga} (17.\,1.\,1882 Wien – 13.\,1.\,1970 Lugano), \emph{Schauspielerin, Sängerin}|pw} und
               bringen Sie mich bei Bubi\pwindex{Schnitzler, Heinrich 9.\,8.\,1902 Hinterbrühl – 12.\,7.\,1982 Wien@\textsc{Schnitzler, Heinrich} (9.\,8.\,1902 Hinterbrühl – 12.\,7.\,1982 Wien), \emph{Regisseur, Schauspieler}|pwv} in
               Erinnerung.\pend
           \pstart Herzlichſt, Ihr \spacefill\mbox{G. Sch.}\pend{}\selectlanguage{ngerman}\endnumbering\briefempfaengerindex{Schnitzler, Arthur@\textsc{Schnitzler, Arthur}!zzzSchwarzkopf, Gustav@\emph{von Gustav Schwarzkopf}!1907-07-171@{17. 7. 1907}|)be}\mylabel{L04271h}
\begin{anhang}
\end{anhang}\newcommand{\dateiname}{L04271}\newcommand{\titel}{Gustav Schwarzkopf an Arthur Schnitzler, 17. 7. 1907}\newcommand{\editorInnen}{Herausgegeben von Jahnke, SelmaMüller, Martin Anton}\input{../tex-inputs/latex-leseansicht-abspann}
